% Source: https://artofproblemsolving.com/community/c3h1569700_modular_arithmetic
% Stable identity: aops-topic-1569700-q01
\textbf{(a) Necessity.} If $x$ is a solution, then $x=b+ck=f+gl$ for some integers $k,l$, so $b-f=gl-ck$ is divisible by $d$.

\textbf{Sufficiency.} Suppose $b-f=dt$. By B\'ezout's identity there are integers $u,v$ with $cu+gv=d$. Put $x=b-cut$. Then $x\equiv b\pmod c$, and
\[
x-f=(b-f)-cut=dt-cut=(cu+gv)t-cut=gvt\equiv0\pmod g .
\]
So $x$ solves both congruences.

\textbf{One residue class.} If $x$ and $y$ are solutions, then $x-y$ is a common multiple of $c$ and $g$, hence a multiple of $L=\operatorname{lcm}(c,g)$; conversely, adding a multiple of $L$ to a solution gives a solution. $\blacksquare$

\textbf{(b)} Here $\gcd(8,9)=1$, so there is exactly one solution modulo $72$. Write $x=3+8k$. Then $3+8k\equiv5\pmod9$, i.e.\ $8k\equiv2\pmod9$. Since $8\equiv-1\pmod9$, this says $-k\equiv2$, so $k\equiv7\pmod9$ and $x=3+8\cdot7=59$. Hence
\[
x\equiv\boxed{59\pmod{72}}.
\]
Check: $59=7\cdot8+3=6\cdot9+5$.

\textbf{(c)} Here $\gcd(6,10)=2$ does not divide $4-1=3$, so by (a) there is $\boxed{\text{no solution}}$. (Directly: the first congruence makes $x$ odd and the second makes $x$ even.)
